\documentclass[11pt]{article}
\pdfinfoomitdate=1\pdftrailerid{}
\usepackage[letterpaper,margin=.7in,headheight=15pt,headsep=18pt,footskip=24pt]{geometry}
\usepackage[T1]{fontenc}
\usepackage{lmodern,amsmath,amssymb,array,booktabs,fancyhdr,tikz}
\renewcommand{\familydefault}{\sfdefault}
\setlength{\parindent}{0pt}\setlength{\parskip}{7pt}
\pagestyle{fancy}\fancyhf{}
\fancyhead[L]{\small Week 87 / Voting and agendas / K--1 adult guide}
\fancyfoot[L]{\scriptsize Bellingham Math Circle / Week 87 / W87-K-FAC-v1}
\fancyfoot[R]{\thepage}
\newcommand{\head}[1]{\par\textbf{\large #1}\par}
\newcommand{\sub}[1]{\par\textbf{#1}\par}
\begin{document}
\head{Mathematical overview: fixed votes, different winners}
\textbf{A majority can go in a circle.} With three fixed ranking cards ABC, BCA, CAB (letters read from first to last), A beats B, B beats C and C beats A, each two votes to one. Each card has a consistent ranking, but the three majority results do not. A child can point to the two supporting cards for each result and say, ``Each one beats someone and loses to someone.''

\textbf{Waiting can decide a knockout, but need not.} Two letters meet first; that winner meets the waiting letter. On these cyclic cards the waiting letter wins. There are only three waiting choices; switching a pair's left and right positions changes no vote. In any three-letter profile with no tied pair result, a letter wins \emph{every} knockout exactly when it beats both others. If it beats both, it wins whether it waits or plays first. If it loses to another letter, meeting that letter first eliminates it. Children can show these two possibilities on the mat, rather than memorize a rule.

\textbf{Being first is sufficient on two cards, not necessary on two.} Putting B first on two of three cards makes B beat both other letters. B can also win every knockout when first on just one card, but never when first on none. In the latter case no card can support B against \emph{both} rivals; two votes against each rival would need more separate cards than the three available. The last page gives a grouping explanation and the complete one-first-card possibilities.

\textbf{Assumptions and limits.} Each of three imaginary voters ranks A, B and C once, with no equal ranks. Each votes for the higher of the \emph{two offered} letters, even when its first choice is absent. All three vote, so a pair cannot tie. The cards stay unchanged throughout a knockout. The child's own favorite does not vote. These finite facts do not establish a fairness claim or Arrow's theorem; changing cards between rounds is a different investigation.

\textbf{What this table reaches.} Problems 1--3 reach the week's cycle and agenda contrast with objects. Problem 4 and the later card-building problems deepen that contrast. Grades 2--5 use the shared card packet; upper readiness adds explanations and profile design. The separate line/median destination belongs to another visit: see the base guide's mathematical overview, page 1. It is not part of this K--1 route.

\textbf{Evidence children can give.} One completed vote or knockout is an experiment. ``The waiting letter always wins'' is initially a conjecture about the current cards, not all cards. Trying all three waiting letters settles that finite question; pointing to the cards that force both pair wins explains why a fixed winner survives every agenda. A few unsuccessful builds do not prove impossibility.

\textbf{Readiness and status.} Children recognize A--F, match V1--V3 to tokens, compare higher/lower and count to three. Independent reading or writing is unnecessary. An anchored adult reads prompts and mat labels, checks each token against its card and records a child's choices when asked; children choose ballots, waiting letters and builds. With two pairs, stagger the first knockout rather than have the adult operate both. Unpiloted; cutting, print scale, slip fit, two-mat handling, resets and timings are physically unrehearsed.

\newpage
\head{Preparation, shared launch and a first visit}
\textbf{Print per pair.} One six-page \texttt{k-1.pdf} (W87-K-v1), plus one copy each of \texttt{week-87-materials.pdf} (W87-M-v1) pages 1 and 4. Add page 2 for Problems 5--8. Print US Letter, single-sided, Actual Size (100\%); check the materials' 100 mm scale. Two K--1 pairs add two copies of each required materials page, beyond the older tables' kits. Materials pages 3 and 5 are not needed.

\textbf{One pair kit.} Materials page 1 supplies three fixed 45 by 80 mm cards, six 28 mm square option tiles (two each A/B/C), and the three 18 mm vote discs V1--V3. Set V4 and V5 aside. Page 4 stays whole: a 170 by 90 mm tally mat with three 54 mm square wells, and a 175 by 100 mm knockout mat. Card votes never use Tie. Page 2 supplies three blank 45 by 80 mm cards and eighteen 21 mm letter slips (six of each letter); only nine slips are needed on the three cards at once. Keep the rest as spares. The roughly 23 mm card rows are tight fits; check real slips before use.

\textbf{Extra preparation.} Allow an additional 15--25 minutes to print, cut and bag two core K--1 kits; add 10--15 minutes for two blank-card/slip sets and 5--10 minutes to rehearse a whole knockout. These are estimates, not measured times. Supply a pencil, scrap paper and optional matching color dots on cards and vote tokens. For the table demonstration make six temporary 28 mm square labels D,D,E,E,F,F from scrap (about two minutes). No week file prints those movable letters; its D/E/F rankings are pictures, not cut cards.

\textbf{After the whole-group launch.} Let partners handle the real pieces. Revisit the non-task DFE, EFD, FED pictures on student page 1. Label the two real tally wells D and E. Ask one child to place V1's token and a partner to check the higher offered letter; children place the remaining tokens. V3 still votes although its first choice F is absent. Point to the already circled E to show how a winner is recorded. The compact reminder below follows the same input, tally and output convention.
\begin{center}
\begin{tabular}{c@{\quad$\longrightarrow$\quad}c@{\quad$\longrightarrow$\quad}c}
Input: V1 / V2 / V3 & Tally for D or E & Circle winner\\[3pt]
$\begin{array}{ccc}D&E&F\\F&F&E\\E&D&D\end{array}$ &
$\begin{array}{c|c}D&E\\\mathrm{V1}&\mathrm{V2, V3}\end{array}$ &
\tikz[baseline=-.6ex]\node[draw,circle,inner sep=3pt]{E};
\end{tabular}
\end{center}
\textbf{Before Problem 2, not before Problem 1.} Use the same pictures and temporary labels: First pair D/E, Waiting F; E wins, then F beats E, both 2--1. Match student page 2's visual and read the mat labels. Use one copy of each letter on the knockout mat and the second copy to label tally wells. Move each loser to Out, move the first winner and waiting tile into Final pair, and move the final winner to Winner. Clear \emph{all three tokens} before the second vote. Reset tiles between knockouts, never the rankings. Remove the temporary D/E/F labels before A/B/C work. This example has one fixed winner, not a cycle, and reveals no task answer.

\textbf{One-hour route.} Allow 5 minutes for movement, 5 for the shared launch, 5 for table handling/recap, 10 for Problem 1, 15 for the knockout recap and Problem 2, 15 for Problem 3, and 5 to share one surprise. Stop earlier if resets still require adult decisions. Problem 4 is a ready pair's continuation; save Problems 5--8 for another visit. Use the small blank frames on pages 5--6 for adult pencil records, not letter slips. The children build on full-size blank cards. Finishing every page is not the goal.

\newpage
\head{Solutions and held hints: the first route}
Adult shorthand such as ABC means first A, second B, third C. Below, page numbers refer to the six-page student packet. Keep solutions out of children's sight.

\sub{Problem 1, page 1: three pair votes}
Circle \textbf{A, B, C}, respectively. A beats B with V1 and V3; B beats C with V1 and V2; C beats A with V2 and V3. Each result is 2--1. Children can show the two supporting cards rather than explain in sentences. After all three results, ask, ``Is there a letter that beat both others?'' There is none: each has a loss.

\textbf{Held hints.} ``Which two letters are on offer?'' ``Which is higher on this voter's card?'' If a child votes for an unavailable first choice, temporarily cover that letter with a small scrap and ask again. Do not discard the voter or turn the ranking upside down.

\sub{Problem 2, page 2: the three printed knockouts}
Keep the fixed cards ABC, BCA, CAB, not the D/E/F example cards.
\begin{center}\renewcommand{\arraystretch}{1.18}
\begin{tabular}{ccccc}
First pair & Waiting & First winner & Final vote & Circle\\\midrule
B, C & A & B & A beats B & A\\
C, A & B & C & B beats C & B\\
A, B & C & A & C beats A & C
\end{tabular}\end{center}
Every vote is 2--1. These are all three waiting choices. Let children notice the pattern; do not announce ``waiting wins'' first. In each case the first winner loses to the letter that waited.

\textbf{Held hints.} ``Who has not played yet?'' ``Which letters reach the final pair?'' If cards or votes change between rounds, return to the unchanged picture and clear only the vote tokens. A natural continuation is to reverse the left/right order of one first pair and see whether any ballot changes; it does not.

\sub{Problem 3, page 3: CBA, BAC, CBA}
Circle \textbf{C only}. C beats A and B, each 2--1, using V1 and V3. B beats A 3--0. With A waiting, C beats B, then C beats A. With B waiting, C beats A, then C beats B. With C waiting, B beats A, then C beats B. Thus C wins all three; A and B never win.

Put away the cut fixed cards and vote directly from the full-size pictures on this page; children need not copy them. The two CBA cards are \emph{two voters}, so both tokens count. On scrap, the adult may record one short row per tried waiting letter and its final winner, leaving the agenda choices to children. This prevents repeating an attempt without noticing.

\textbf{Held hints.} ``Who is voting on this page?'' ``Can either other letter beat C?'' A child can point to the two cards with C first: neither rival can remove C, whether C waits or plays first. This breaks a universal ``waiting wins'' conjecture without changing the voting rule.

\sub{Problem 4, page 4: ABC, BAC, CAB}
Circle \textbf{A only}. A beats B (V1/V3) and C (V1/V2), each 2--1; B beats C (V1/V2), also 2--1. Waiting A gives B then A; waiting B gives A then A; waiting C gives A then A. Only A wins. Each letter is first on one card, so counting first choices alone would miss the result. Again vote from the printed cards, not the old cut set.

\textbf{Held hints.} ``If A plays first, who can knock it out?'' ``If A waits, who could beat it at the end?'' After testing, ask how these cards differ from Problem 1: only V2 has exchanged A and C. That comparison prepares later card changing without solving Problem 5's different target.

\newpage
\head{Later visit: change and build the voters}
\sub{Problem 5, page 5: one swap to make C win}
Return to the pictured original cards ABC, BCA, CAB. Exactly three swaps work: \textbf{V1 ABC to ACB} (swap B/C); \textbf{V1 ABC to CBA} (swap A/C); or \textbf{V2 BCA to CBA} (swap B/C). C beats B 2--1 in all three; C beats A 2--1, 3--0 and 2--1, respectively. Thus C wins every knockout. All other cards stay fixed. Make only the replacement card with slips and put the original aside. Reset to the original three cards before testing another swap.

This list is complete: each of three cards offers three pairs of positions to exchange; checking those nine swaps gives just the three above. Children need one working construction, not this adult catalog. \textbf{Held hints:} ``Which rival stops C on the original cards?'' ``Which card could change its vote between B and C?'' C already beats A; the changed ranking must also let C beat B without losing that advantage.

\sub{Problem 6, page 5: build a fixed winner B}
One example is \textbf{BAC, BCA, ABC}: B beats A 2--1 and C 3--0, and wins all three knockouts. Repeated cards are allowed; the older shared packet's ``three different cards'' rule does not apply here. Any legal set with B first on two or three cards works, since those cards support B against either rival. Accept a different build after testing its two pair votes and explaining the agendas; putting B first everywhere is a valid entry, not a failure of imagination.

\textbf{Held hints.} ``Which two cards could keep voting for B against either other letter?'' ``Does your set work when B plays first as well as when B waits?'' Let the child arrange the slips; the adult records the final ranking in the small frames. To deepen a quick success, move to the next problem rather than demanding many similar builds.

\sub{Problem 7, page 6: B first exactly once}
\textbf{It can be done.} The only working sets, ignoring which voter receives which card, are
\[
\{ABC,\ BAC,\ CBA\}\qquad\text{and}\qquad\{ABC,\ BCA,\ CBA\}.
\]
In either, B beats A and C, both 2--1. Assign the three cards to V1--V3 in any order. The B-first card supplies one vote against each rival. To get a second vote in both comparisons, the other two cards must put B second, one below A (ABC) and one below C (CBA). A B-last card supplies neither vote. The B-first card may be BAC or BCA; these exhaust the possibilities.

\textbf{Held hints.} ``Which of the two rivals still beats B?'' ``What can a card with B second help B beat?'' Do not hand over the answer after a few random failures. Children can keep their one B-first card and rearrange the others. A ready pair can find both working sets; this is optional, not a new written task.

\sub{Problem 8, page 6: B never first}
\textbf{Impossible.} Draw two separate loops on scrap, labelled ``B above A'' and ``B above C.'' Sort each voter's token by its card: B second fits one loop; B last fits neither. With B never first, no token qualifies for both. For B to beat both rivals, each loop would need two different voter tokens, but only three exist. This is a grouping proof, not two simultaneous votes. Children may simply point to the shortage; no written counting argument is required.

\textbf{Held hints.} ``Can a card with B not first put B above both A and C?'' ``Can both loops have two of our three tokens without using one twice?'' Restore each token to its card before ordinary voting. Failed builds alone are not the explanation; the shortage applies to every legal build.

\textbf{Where it goes.} Once the pages are approved, this section replaces ``K--1 route (adult-led)'' on page 1 of \texttt{week-87-facilitator.pdf} (W87-F-v2). Its preparation adds the K--1 kits to that guide's older-table plan; the older-band overview and solutions remain in place.
\end{document}
